\section{A solvable comparison model: flattening can preserve kernel ordering}
\label{sec:linear-control-comparison}

The following auxiliary model isolates an obstruction to inferring
posterior behavior from the preceding flattening theorem. It is a
Gaussian covariance-control model, not an identification of the tanh
network's critical rate.

Fix a positive-definite matrix $K$ on a finite input set and an integer
$H\ge0$. For a nonzero output vector $z$, define
\[
 I_{H,K}(z)=\inf_{Q_1,\ldots,Q_H\succ0}
 \left\{\sum_{j=1}^H D\bigl(\mathcal N(0,Q_j)
                  \Vert\mathcal N(0,Q_{j-1})\bigr)
           +\frac12z^TQ_H^{-1}z\right\},\qquad Q_0=K.
\]
There is no tanh transition in this definition.

\begin{proposition}[Exact Gaussian covariance-control value]
\label{prop:linear-control-value}
Put $s=z^TK^{-1}z$. For $H\ge1$, let $a>1$ be the unique solution of
\[
 s=a^H(a-1).
\]
Then
\[
 I_{H,K}(z)=\frac12\{H(a-1-\log a)+(a-1)\}.
\]
For $H=0$ the value is $s/2$. For every fixed $H$ the value is strictly
increasing in $s$, while for every fixed $s>0$ it tends to zero as
$H\to\infty$.
\end{proposition}

\begin{proof}
Apply the invertible transformation $K^{-1/2}$ to all Gaussian vectors
and to $z$. Relative entropy is unchanged, so the initial covariance
becomes the identity. Write $u=K^{-1/2}z/\sqrt{s}$ and
$q_j=u^T(K^{-1/2}Q_jK^{-1/2})u$, with $q_0=1$.
Data processing under projection onto $u$ gives
\[
 D\bigl(\mathcal N(0,Q_j)\Vert\mathcal N(0,Q_{j-1})\bigr)
 \ge\frac12\left\{\frac{q_j}{q_{j-1}}-1
                     -\log\frac{q_j}{q_{j-1}}\right\}.
\]
Cauchy--Schwarz gives $z^TQ_H^{-1}z\ge s/q_H$.
Both inequalities are attained simultaneously by keeping the covariance
equal to the identity on $u^\perp$ and changing only its variance on
$u$. Hence, putting $b_j=\log(q_j/q_{j-1})$, the exact value is
\[
 \inf_{b_1,\ldots,b_H\in\mathbb R}
 \frac12\left\{\sum_{j=1}^H(e^{b_j}-1-b_j)
                    +s\exp\left(-\sum_{j=1}^H b_j\right)\right\}.
\]
At a fixed sum of the $b_j$, Jensen's inequality makes them all equal.
The remaining strictly convex scalar objective has derivative zero
exactly when $a-1=s a^{-H}$. Substitution gives the formula.
Strict monotonicity follows by differentiating the minimized objective
with respect to $s$, giving derivative $1/(2a^H)>0$.
For convergence to zero, take every $b_j=2\log H/H$ for $H\ge2$.
Then the transition cost is $O((\log H)^2/H)$ and the readout cost is
$s/(2H^2)$.
\end{proof}

For a prescribed sign pattern $f$ and margin $r>0$, minimize this value
over $f(x)z(x)\ge r$. Its minimizers are exactly the minimizers of
$z^TK^{-1}z$ under those constraints, at every $H$. The same is true
when only training-point sign constraints are imposed. Thus this
comparison model has rates that flatten to a common value while its
minimum-rate sign predictions retain the kernel minimum-norm ordering
at every finite depth.

This demonstrates why absolute flattening alone cannot answer the
target-relative question for the actual tanh learner. To use this model
as more than a comparison one would have to control the nonlinearity
and arbitrary non-Gaussian feature-law controls at the scale of the
remaining rate differences. No such comparison is established here.
